\chapter{案例：一阶 bang-bang 系统与等效控制}\label{ch:case}

前两章给出两套理论：右端 Lipschitz 时使用经典解（第\ref{ch:lipschitz}章），右端有界时使用 Filippov 解与集值工具（第\ref{ch:bounded}章）。本章把它们落到两个可以完全手算的一阶系统上：
\begin{equation}\label{eq:case-system}
  \text{(S1)}\quad \dot x=-\sgn{x},
  \qquad\qquad
  \text{(S2)}\quad \dot x=-\sgn{x}+0.9\sin t .
\end{equation}
两条系统的控制幅值都是 $k=1$，只差一个时变扰动 $0.9\sin t$，而 $0.9<1$ 恰好落在控制幅值以内。这个“差一点”的扰动足以让三种解概念、三种证明方法出现分歧，因此它们能说明本章要讲的全部要点：
\begin{enumerate}
  \item 三种解概念（classical / Carathéodory / Filippov）各自能说什么、不能说什么（第\ref{sec:case-solutions}节）；
  \item \textbf{等效控制}（equivalent control）是什么、为什么教材几乎只讲它、它成立需要什么条件（第\ref{sec:case-ec}节，本章重点）；
  \item 两个 Lyapunov 函数 $V_1=\tfrac12x^2$ 与 $V_2=|x|$ 各自能给出什么（第\ref{sec:case-lyapunov}节）；
  \item 三种处理非光滑性的方法：几乎处处导数、Dini 上导数、Clarke 广义梯度与集值 Lie 导数（第\ref{sec:case-methods}节）。
\end{enumerate}

一句话概括全章：\textbf{等效控制负责“滑模面上的运动”，非光滑 Lyapunov 工具负责“怎样到达滑模面”，把两者拼起来才是完整证明。}

\section{两个系统与三种解概念}\label{sec:case-solutions}

两条系统都可写成仿射形式
\[
\dot x=f_i(t,x)+g\,u,\qquad u=-\sgn{x},\qquad s(x)=x,\qquad g=1,
\]
其中 $f_1\equiv0$，$f_2(t,x)=0.9\sin t$。记
\begin{equation}\label{eq:case-fields}
  \mathbb{K}_i:=\mathbb{K}[f_i]\quad(i=1,2),
\end{equation}
即 $f_1=-\sgn{x}$、$f_2=-\sgn{x}+0.9\sin t$ 的 Filippov 集值化。

\subsection{classical 解：两条系统都没有}

先看“吸引”这一半。$x>0$ 时
\[
\text{(S1): }\ \dot x=-1,
\qquad
\text{(S2): }\ \dot x=-1+0.9\sin t\le-0.1,
\]
$x<0$ 时两式符号相反。因此任一初值都在有限时间内到达 $x=0$：$(S1)$ 用时 $|x_0|$，$(S2)$ 至多 $10|x_0|$（估计见第\ref{sec:case-lyapunov}节）。

问题出在到达之后。设轨迹在 $t_0$ 时刻到达 $x=0$。若它要成为 $C^1$ 解，则 $\dot x(t_0)$ 必须同时满足两侧的要求：

\begin{itemize}
  \item \textbf{若穿过原点：} 导数必须在 $t_0$ 连续。左极限是 $-1+\delta(t_0)$（$\delta$ 表示扰动：$(S1)$ 为 $0$，$(S2)$ 为 $0.9\sin t$），而在 $x=0$ 处方程给出的值是 $\delta(t_0)$，二者相差恰为 $1=k$——这个差来自控制量的不连续跳变 $2k$，与扰动无关。
  \item \textbf{若停在原点：} 右导数为 $0$，而左极限仍是 $-1+\delta(t_0)$，二者相差 $-1+\delta(t_0)$，在 $|\delta|\le0.9$ 下其绝对值不小于 $0.1$，同样不可能为零。
\end{itemize}

于是两条系统在 $[0,\infty)$ 上\textbf{都不存在 classical 解}。

$(S1)$ 的最大解 $x(t)=\max\{x_0-t,0\}$ 只是\textbf{分段 classical}（绝对连续）的；它的导数
\[
\dot x(t)=\begin{cases}
  -1,& t<x_0,\\
  \ \ 0,& t>x_0
\end{cases}
\]
在 $t=x_0$ 处不存在。另需注意：$(S1)$ 之所以“看起来 classical 够用”，是因为约定 $\sgn{0}=0$ 恰好使 $0$ 落在 $\mathbb{K}_1(0)=[-1,1]$ 之中，这个取值正好对应“停留在原点”的那条滑动段——这是巧合，而不是经典理论在这里成立。

\subsection{Carathéodory 解：一个巧合与一次失败}

Carathéodory 解只要求绝对连续并在\textbf{几乎处处}意义下满足方程（式~\eqref{eq:caratheodory}）。它对两条系统的判别完全不同。

\textbf{(S1)：恰好有解。} 在约定 $\sgn{0}=0$ 下，常值 $x(t)\equiv0$ 满足 $0=-\sgn{0}$，因此它是 $(S1)$ 的 Carathéodory 解；$x(t)=\max\{x_0-t,0\}$ 也是。

\textbf{(S2)：连全局解都没有。} 设 $x(t)$ 是 $(S2)$ 在 $[0,T)$ 上的绝对连续解，取 $V_2=|x|$。在 $x\neq0$ 处
\[
\dot V_2=\sgn{x}\,\dot x=-1+0.9\sgn{x}\sin t\le-0.1 .
\]
于是 $V_2$ 每秒至少下降 $0.1$，而 $V_2\ge0$，故 $V_2(t)=0$ 对一切 $t\ge10|x_0|$ 成立（积分论证见引理~\ref{lem:case-finite-time}）。此后 $x\equiv0$，而 $(S2)$ 在 $x=0$ 处要求 $\dot x=0.9\sin t$，对几乎所有 $t$ 不为零，矛盾。因此 $(S2)$ \textbf{不存在全局 Carathéodory 解}。

教训是：Carathéodory 的“几乎处处”只放松了\textbf{时间}方向的光滑性，右端仍必须是单点等式；一旦方程在切换面上与“停留”直接矛盾（$0=0.9\sin t$），再放松时间方向也救不回来。

\subsection{Filippov 解：把矛盾点换成集合}

Filippov 的做法是承认这个矛盾，并把 $x=0$ 处的取值换成集合（定义~\ref{def:filippov-solution}、式~\eqref{eq:filippov-set}）。对两条系统：
\begin{equation}\label{eq:case-filippov}
  \mathbb{K}_1(0)=\overline{\mathrm{co}}\{-1,\,1\}=[-1,\,1],
  \qquad
  \mathbb{K}_2(t,0)=0.9\sin t+[-1,\,1]=[-1+0.9\sin t,\;1+0.9\sin t].
\end{equation}
两个集合都包含 $0$；并且由于 $|0.9\sin t|\le0.9<1$，$0$ 是\textbf{内点}。于是 $x\equiv0$ 同时是两条系统的 Filippov 解，$0\in\mathbb{K}_i$ 正是“允许停留”的精确含义。

更强地，$\{0\}$ 是\textbf{强不变}的：任何解都在有限时间内到达原点，而到达之后不能离开——若在某个区间上 $|x|>0$，则该区间上 $|x|$ 严格下降，与非负性（起点为 $0$）矛盾。注意这需要两件事同时成立：$0\in\mathbb{K}_i$（允许停留）与两侧向量场都指向滑模面（$|x|$ 在 $x\neq0$ 时严格下降）。缺任何一条，原点都不是吸引的。

\begin{table}[!htb]
  \centering
  \small
  \begin{tabular}{@{}p{2.2cm}|p{3.2cm}|p{2.2cm}|p{3.4cm}|p{2.8cm}@{}}
    \hline
    解概念 & 对右端的要求 & 光滑性 & (S1) & (S2) \\
    \hline
    classical & 关于 $x$ 连续（唯一性需 Lipschitz） & $C^1$ & 无全局解（导数在到达时刻跳变） & 无全局解 \\
    \hline
    Carathéodory & 关于 $x$ 连续，关于 $t$ 可测 & 绝对连续 & 有解（靠 $\sgn{0}=0$ 的巧合） & \textbf{无全局解} \\
    \hline
    Filippov & 可测 + 局部本质有界 & 绝对连续 & 有解，$x\equiv0$ 且 $0$ 强不变 & 有解，$x\equiv0$ 且 $0$ 强不变 \\
    \hline
  \end{tabular}
  \caption{三种解概念在两条系统上的对照}
  \label{tab:case-solution-concepts}
\end{table}

顺带把三条存在性定理的假设对齐，可以看出本例落在哪个缝隙里：

\begin{itemize}
  \item classical 解的\textbf{局部存在性}只要求右端连续（Peano 定理），\textbf{唯一性}才要求局部 Lipschitz（第\ref{ch:lipschitz}章定理~\ref{thm:local-existence}）；
  \item Carathéodory 定理要求右端关于 $x$ 连续，$(S2)$ 在 $x=0$ 处不满足；
  \item Filippov 解只要求右端可测且局部本质有界（定理~\ref{thm:filippov-existence}）。
\end{itemize}

$(S1)$、$(S2)$ 的右端既不连续、也不局部 Lipschitz，前两条都用不上；而它们可测且局部有界，Filippov 存在性定理恰好在“右端不连续”这条缝隙里生效。

\section{等效控制：Utkin 范式}\label{sec:case-ec}

上一节说明：在非光滑框架里，滑模面上的“速度”是一个集合。工程上真正被广泛使用的描述方式却不是微分包含，而是 Utkin 的\textbf{等效控制}：把切换信号换成一个等效的连续信号，在滑模面上做常规（光滑）分析。本节先给出定义，再在 $(S1)$、$(S2)$ 上手算，然后说明它的严格含义、成立条件，以及它能回答什么、不能回答什么。

\subsection{定义：把切换信号换成连续信号}

考虑
\begin{equation}\label{eq:case-affine}
  \dot x=f(t,x)+g(t,x)u,\qquad s=s(t,x),\qquad u=-k\,\sgn{s},\quad k>0 .
\end{equation}
沿轨迹对 $s$ 求导：
\[
\dot s=\partial_t s+\nabla s^\top\big(f+gu\big)
=\underbrace{\big(\partial_t s+\nabla s^\top f\big)}_{\text{漂移项}}
+\underbrace{\big(\nabla s^\top g\big)}_{\text{控制增益}}\,u .
\]
若 $\nabla s^\top g\neq0$，令 $\dot s=0$ 解出
\begin{equation}\label{eq:case-ueq}
  u_\mathrm{eq}(t,x)=-\frac{\partial_t s+\nabla s^\top f}{\nabla s^\top g},
\end{equation}
这就是 Utkin 的\textbf{等效控制}。把它代回系统，得到\textbf{降维等效动力学}
\[
\dot x=f+g\,u_\mathrm{eq},
\]
其中不再含不连续项，可以按连续系统做分析。教材里“滑模阶段 $\dot s=0$”那句话，说的就是这个操作。

需要立刻指出：这一步在形式上是“先假设 $s$ 可微、再令其导数为零”，而由 \S\ref{sec:case-solutions} 已经看到 $s(t)=x(t)$ 在切换时刻的导数会跳变 $1$。所以式~\eqref{eq:case-ueq} 必须另有严格解释，见 \S\ref{sec:case-ec-meaning}。

\subsection{手算：(S1) 与 (S2) 的等效控制}

对两条系统，$s=x$、$\partial_t s=0$、$\nabla s=1$、$g=1$，于是 $u_\mathrm{eq}=-f$，即“等效控制恰好抵消漂移项”：

\begin{center}
\small
\begin{tabular}{@{}c|c|c|c|c@{}}
  \hline
  系统 & 漂移 $f(t,x)$ & 等效控制 $u_\mathrm{eq}$ & 滑动方程 $\dot x=f+gu_\mathrm{eq}$ & 可容许性 \\
  \hline
  (S1) & $0$ & $0$ & $0$ & $|u_\mathrm{eq}|=0<1$ \\
  \hline
  (S2) & $0.9\sin t$ & $-0.9\sin t$ & $0.9\sin t-0.9\sin t=0$ & $|u_\mathrm{eq}|\le0.9<1$ \\
  \hline
\end{tabular}
\end{center}

三点观察。

\begin{enumerate}
  \item \textbf{滑模面上的运动与扰动无关。} $(S2)$ 的等效控制是时变的，恰好把扰动“抵消”掉，因此降维后的动力学与无扰动的 $(S1)$ 完全相同（都是 $\dot x=0$）。这就是滑模控制最核心的性质：只要扰动与控制进入同一个通道（匹配条件，$g=1$ 对两条系统都成立）且幅值小于 $k$，滑模面上的运动就完全不受它影响。
  \item \textbf{等效控制不必是常值。} 它可以是时变的、甚至状态相关的，只要落在可容许区间内（第\ref{sec:case-ec-conditions}节）。
  \item \textbf{差别只体现在到达过程。} 两条系统的滑动方程相同，但 $(S2)$ 的到达时间上界放宽了 $10$ 倍（$10|x_0|$ 对 $|x_0|$）——扰动被抑制在滑模面上，代价体现在趋近段。
\end{enumerate}

\subsection{等效控制到底等于什么：Filippov 集的法向零点}\label{sec:case-ec-meaning}

式~\eqref{eq:case-ueq} 的严格含义可以用一句话说清：\textbf{等效控制是 Filippov 集里让滑模面法向速度为零的那个（唯一）控制值}。把它写成命题。

\begin{proposition}[等效控制 = 唯一的法向零点]\label{prop:case-ueq}
  考虑~\eqref{eq:case-affine}，设 $s=s(x)$ 与 $t$ 无关，$g$ 与 $u$ 无关，切换值在 $s\to0^\pm$ 时的极限为 $u^\pm$。若在滑模面 $s=0$ 上
  \[
  \nabla s(x)^\top g(x)\neq0,
  \qquad
  \big(\nabla s^\top(f+gu^-)\big)\big(\nabla s^\top(f+gu^+)\big)<0 ,
  \]
  则
  \[
  \big\{v\in\mathbb{K}[f+gu](x):\ \nabla s(x)^\top v=0\big\}=\big\{f+g\,u_\mathrm{eq}\big\},
  \]
  并且相应的混合值 $u_\mathrm{eq}$ 落在 $\mathrm{co}\{u^-,u^+\}$ 内部。
\end{proposition}

\begin{proof}
  在滑模面上 $\mathbb{K}[f+gu]=\mathrm{co}\{f+gu^-,f+gu^+\}=f+g\,\mathrm{co}\{u^-,u^+\}$。
  泛函 $u\mapsto\nabla s^\top(f+gu)$ 是斜率为 $\nabla s^\top g\neq0$ 的仿射函数，故严格单调；两端取值异号，故在 $\mathrm{co}\{u^-,u^+\}$ 内有唯一零点 $u_\mathrm{eq}$，且 $\nabla s^\top g\neq0$ 使 $u_\mathrm{eq}=-(\nabla s^\top f)/(\nabla s^\top g)$。$v$ 与 $u$ 一一对应，结论得证。
\end{proof}

命题~\ref{prop:case-ueq} 解释了三件事。

\begin{enumerate}
  \item \textbf{为什么教材可以不理 Filippov。} 在“好条件”（横截条件 + 滑动条件）下，Filippov 集与切平面 $\{\nabla s^\top v=0\}$ 的交是单点，而这个点恰好就是 $f+gu_\mathrm{eq}$。两条路线给出同一个速度，因此用等效控制做的分析在数学上是成立的。
  \item \textbf{“抖振的平均”是什么。} 取 $u=-k\,\sgn{s}$，设滑模面上控制取 $+k$ 的时间比例为 $\lambda_+$，则平均控制为 $k\lambda_+-k(1-\lambda_+)=k(2\lambda_+-1)$，令其等于 $u_\mathrm{eq}$ 得
  \[
  \lambda_+(t)=\frac{k+u_\mathrm{eq}(t)}{2k}\in(0,1).
  \]
  对 $(S2)$：$\lambda_+(t)=(1+0.9\sin t)/2\in[0.05,\,0.95]$。物理图像很直白：扰动为正时略多花一点时间取 $+1$ 去抵消它，扰动为负时反过来。
  \item \textbf{为什么 $0\in\partial\mathbb{K}$ 时滑模会退化。} 命题要求 $u_\mathrm{eq}$ 在 $\mathrm{co}\{u^-,u^+\}$ \textbf{内部}（这正是 $0\in\operatorname{int}\mathbb{K}$）。若 $u_\mathrm{eq}$ 恰好落在端点，法向分量没有严格的“两侧异号”，滑模只剩擦过（grazing）的情形。
\end{enumerate}

\subsection{等效控制成立的四个条件}\label{sec:case-ec-conditions}

式~\eqref{eq:case-ueq} 的推导只用了一行代数，但要让结论有意义，需要以下条件\emph{同时}成立：

\begin{enumerate}
  \item \textbf{仿射结构：} 控制以加性方式进入系统，$\dot x=f(t,x)+g(t,x)u$；
  \item \textbf{单张光滑滑模面 + 横截条件：} $s$ 光滑单值，$\{s=0\}$ 是一张光滑超曲面，且 $\nabla s^\top g\neq0$ 在其上处处成立（保证 $u$ 能影响 $\dot s$）；
  \item \textbf{滑动（可容许）条件：} $u_\mathrm{eq}(t,x)\in\mathrm{co}\{u^-,u^+\}$，即 $|u_\mathrm{eq}|\le k$，等价于 $0\in\operatorname{int}\mathbb{K}[f+gu]$；
  \item \textbf{可达性：} 存在控制使状态在有限时间内到达 $s=0$（趋近段 $V=\tfrac12s^2$ 严格下降）。
\end{enumerate}

条件 3 最容易被忽略，而它恰恰是工程上最本质的一条：\textbf{等效控制的幅值不能超过控制幅值}。$(S2)$ 中的 $0.9$ 正好是这条界的临界参数。把扰动幅度改成 $1.5$：
\[
\dot x=-\sgn{x}+1.5\sin t,\qquad u_\mathrm{eq}(t)=-1.5\sin t .
\]
当 $|\sin t|>2/3$ 时 $|u_\mathrm{eq}|>1$，不再可容许。此时
\[
\mathbb{K}(t,0)=[-1+1.5\sin t,\;1+1.5\sin t]
\]
不含 $0$：原点不再是解，轨迹被推出滑模面；等到 $|\sin t|\le2/3$ 时它又被拉回。这就是\textbf{间歇滑模}（intermittent sliding）——滑模时有时无，等效控制时而存在、时而“欠幅”。

用 Khalil 的语言说，条件 3 就是：匹配扰动的幅值必须小于控制增益，$|f|/|g|<k$ \cite[§14.1]{khalilNonlinearSystems}。这不是技术性假设，而是“控制器是否买得起抗扰性能”的预算约束。

\subsection{等效控制能给出什么、不能给出什么}

\begin{itemize}
  \item \textbf{能给出：} 滑模面上的不变性（$x\equiv0$ 是解）与降维动力学；等效控制幅值的需求（据此选 $k$）；工程实现所需的占空比 $\lambda_+$。
  \item \textbf{不能给出：} 到达时间与到达前的瞬态（那是滑模面之外的世界，需要趋近段的 Lyapunov 分析，见第\ref{sec:case-lyapunov}节）；条件 3 失效时的行为（$u_\mathrm{eq}$ 不存在，只能回到集值描述）；横截条件破坏的奇异情形（第\ref{sec:utkin-counterexample}节的反例）。
  \item \textbf{尤其不能：} 用 $u_\mathrm{eq}$ 直接“证明”稳定性。$u_\mathrm{eq}$ 是“假设滑模已经发生”之后的产物，若没有独立证明滑模确实发生（条件 2--4），用降维系统得到的稳定性结论是循环论证。这正是本章要引入非光滑工具的原因。
\end{itemize}

\subsection{为什么教材普遍只用等效控制}

绝大多数滑模教材采用“趋近阶段 + 滑动阶段”的分工，因而可以完全绕开非光滑工具：

\begin{enumerate}
  \item \textbf{趋近阶段（$s\neq0$）：} $\sgn{s}$ 是常值，右端光滑，普通 Lie 导数完全够用。取 $V=\tfrac12s^2$，只需验证 $\dot V=s\dot s<0$ 即可证明可达性与到达时间。
  \item \textbf{滑动阶段（$s=0$）：} 用等效控制把不连续系统换成降维的连续系统 $\dot x=f+gu_\mathrm{eq}$，后续分析与连续系统无异（本节的 $(S1)$、$(S2)$ 都化成 $\dot x=0$）。
  \item \textbf{Lyapunov 函数：} 选用光滑的 $V=\tfrac12s^2$，其梯度 $\nabla V=s$ 处处存在，天然避开 Clarke 广义梯度。
\end{enumerate}

如果确实需要非光滑 Lyapunov 函数（例如 $V=|s|$），教材与论文中更常用的是 \textbf{Dini 上导数}：它的定义只依赖函数本身与单点向量场，不需要引入广义梯度，配套的 Lyapunov 定理（定理~\ref{thm:dini-lyapunov}）与经典 LaSalle 不变集原理（\cite{roucheStabilityTheoryLiapunov1977}）都有现成版本。交集型\textbf{集值 Lie 导数}则是专为 Filippov 版 LaSalle 不变集原理（定理~\ref{thm:LaSalle}）设计的工具，属于理论论文分支，因此没有成为教材标配——第\ref{sec:case-methods}节会看到，它在原点处给出的信息其实相当有限。

简而言之：教材教的是怎么“用”滑模控制（等效控制 + 光滑 $V$），论文讨论的是滑模控制的“数学基础”（Filippov 微分包含 + 集值李导数）。二者分工不同，并不矛盾。

\section{两个 Lyapunov 函数}\label{sec:case-lyapunov}

有了等效控制，滑模面上的问题已经解决；剩下的问题是“怎样到达滑模面”以及“到达之后会不会离开”。后者由 \S\ref{sec:case-solutions} 的强不变性回答，前者需要 Lyapunov 函数。本节比较两个最常用的选择：光滑的 $V_1=\tfrac12x^2$ 与非光滑的 $V_2=|x|$。

\subsection{经典 Lyapunov 定理与它的两个前提}

作为对照，先写出经典（光滑）版本 \cite[Theorem 4.1]{khalilNonlinearSystems}。考虑自治系统 $\dot x=f(x)$，$f$ 在包含原点的区域 $D$ 上局部 Lipschitz，$f(0)=0$。若存在连续可微的 $V:D\to\RR$ 满足 $V(0)=0$、$V(x)>0$（$x\in D$，$x\neq0$）且
\[
\dot V(x)=\nabla V(x)^\top f(x)\le0\qquad(x\in D),
\]
则原点稳定；若 $\dot V(x)<0$（$x\neq0$），则原点渐近稳定。这个定理的每一字都依赖两个隐含前提：

\begin{enumerate}
  \item \textbf{右端映射单点：} $f(x)$ 在每一点都有唯一确定的值；
  \item \textbf{$V$ 光滑可微：} 梯度 $\nabla V(x)$ 是单点。
\end{enumerate}

两条系统恰好把这两个前提同时打破（速度是集合 $\mathbb{K}_i$，$V_2$ 在原点不可微），因此经典定理\textbf{不能}直接用于 $(S1)$、$(S2)$——即便 $V_1$ 是光滑的，因为系统本身没有 classical 解，定理关于解的存在性前提已经落空。这就是必须先换解概念、再谈 Lyapunov 工具的原因。

\subsection{$V_1=\tfrac12x^2$：光滑，但在原点处导数为零}

$V_1$ 光滑、正定、径向无界。沿轨迹（$x\neq0$ 时）：
\[
\text{(S1): }\ \dot V_1=x\dot x=-|x|=-\sqrt{2V_1},
\]
\[
\text{(S2): }\ \dot V_1=x\big(-\sgn{x}+0.9\sin t\big)=-|x|+0.9\,x\sin t\le-0.1|x|=-0.1\sqrt{2V_1}.
\]
两式都是 $\dot V_1\le-\alpha\sqrt{2V_1}$ 的形式（$(S1)$ 取 $\alpha=1$，$(S2)$ 取 $\alpha=0.1$）。由比较引理~\ref{lem:comparison}（或例~\ref{ex:reaching-time} 的推导），有限时间到达估计为
\[
t_r\le\frac{\sqrt{2V_1(0)}}{\alpha}=\frac{|x_0|}{\alpha},
\]
即
\begin{equation}\label{eq:case-bounds}
  \text{(S1): }\ T_1\le|x_0|,
  \qquad\qquad
  \text{(S2): }\ T_2\le10\,|x_0| .
\end{equation}
对 $(S1)$，界是精确的（$x(t)=\max\{x_0-t,0\}$）。

要注意 $V_1$ 的局限：它在原点处 $\dot V_1=0$，所以“严格下降”只在 $\{x\neq0\}$ 上成立。渐近性不能由 $\dot V_1<0$ 直接得到，而要靠“有限时间内到达 + 原点强不变”。另外，$\dot V_1\le-\alpha\sqrt{2V_1}$ 这种“$\sqrt{V}$ 型”不等式给出的才是有限时间，若只有 $\dot V\le-\alpha V$ 则只能得到指数收敛（见例~\ref{ex:reaching-time}）。

\subsection{$V_2=|x|$：非光滑，但下降速率有下界}

$V_2$ 正定、径向无界，但在原点不可微。沿轨迹（$x\neq0$ 时）：
\[
\text{(S1): }\ \dot V_2=\sgn{x}\,\dot x=-\sgn{x}\,\sgn{x}=-1,
\]
\[
\text{(S2): }\ \dot V_2=\sgn{x}\big(-\sgn{x}+0.9\sin t\big)=-1+0.9\,\sgn{x}\sin t\le-0.1 .
\]
与 $V_1$ 相比，$V_2$ 给出的信息更强：下降速率有与状态无关的正下界。因此
\[
\text{(S1): }\ T_1\le\frac{|x_0|}{1}=|x_0|,
\qquad\qquad
\text{(S2): }\ T_2\le\frac{|x_0|}{0.1}=10\,|x_0|,
\]
与式~\eqref{eq:case-bounds} 一致——这不是巧合，因为 $\sqrt{2V_1}=V_2$。

代价是 $V_2$ 在原点不可微，普通 Lie 导数在 $x=0$ 处无定义。第\ref{sec:case-methods}节的三种方法正是为了合法地在原点附近谈论 $\dot V_2$。

\begin{table}[!htb]
  \centering
  \small
  \begin{tabular}{@{}p{2.6cm}|p{4.6cm}|p{4.6cm}@{}}
    \hline
     & $V_1=\tfrac12x^2$ & $V_2=|x|$ \\
    \hline
    光滑性 & 处处 $C^1$ & 原点不可微 \\
    \hline
    下降速率 & $\dot V_1\le-\alpha\sqrt{2V_1}$（依赖状态） & $\dot V_2\le-\alpha$（常数，更强） \\
    \hline
    原点处 & $\dot V_1=0$，需另证有限时间 & 普通导数无定义，需 a.e./Dini/Clarke \\
    \hline
    能否直接用经典定理 & 系统无 classical 解，不能 & 不能（非光滑且非 classical） \\
    \hline
    主要用途 & 趋近阶段可达性（教材范式） & 有限时间收敛的严格证明 \\
    \hline
  \end{tabular}
  \caption{两个 Lyapunov 函数的对比}
  \label{tab:case-lyapunov}
\end{table}

\section{三种方法：几乎处处导数、Dini 上导数、Clarke 广义梯度}\label{sec:case-methods}

三种方法的差别不在“用什么函数”，而在“怎样在原点附近合法地谈论导数”。首先给出三种方法都要用到的一个引理。

\begin{lemma}[非负绝对连续函数的有限时间归零]\label{lem:case-finite-time}
  设 $t\mapsto V(t)\ge0$ 绝对连续，且存在 $c>0$ 使 $\dot V(t)\le-c$ 在集合 $\{V>0\}$ 上几乎处处成立。则
  \[
  V(t)=0\qquad\text{对一切 } t\ge V(0)/c .
  \]
\end{lemma}

\begin{proof}
  在 $\{V>0\}$ 的每个连通区间上积分得 $V$ 以速率 $c$ 下降，故 $V$ 必在 $V(0)/c$ 之前归零。另一方面 $V$ 不能从 $0$ 重新变为正：若 $V(t_0)=0$ 而 $V>0$ 于 $(t_0,t_1]$，则在该区间上积分给出 $V(t_1)\le V(t_0)-c(t_1-t_0)<0$，与 $V\ge0$ 矛盾。
\end{proof}

引理~\ref{lem:case-finite-time} 同时给出了“有限时间到达”与“到达后不离开”两件事，因此它是三种方法共同的落脚点。

\subsection{方法一：几乎处处导数（推荐）}

绝对连续函数的复合仍是绝对连续的：$V_2$ 局部 Lipschitz、$x(\cdot)$ 绝对连续，故 $V_2(x(\cdot))$ 绝对连续，并且几乎处处成立链式法则
\[
\frac{\dd}{\dd t}V_2(x(t))=\sgn{x(t)}\,\dot x(t)\qquad\text{a.e.}
\]
在 $x\neq0$ 处上式给出
\[
\text{(S1): }\ \dot V_2=-1,
\qquad
\text{(S2): }\ \dot V_2=-1+0.9\,\sgn{x}\sin t\le-0.1 ,
\]
于是引理~\ref{lem:case-finite-time} 直接给出式~\eqref{eq:case-bounds} 的界，并且同时保证“到达后不离开”。结论：

\begin{quote}
$(S1)$、$(S2)$ 在 Filippov 意义下\textbf{全局有限时间稳定}，到达时间界分别为 $|x_0|$ 与 $10|x_0|$；$(S1)$ 的界精确。
\end{quote}

这条证明只有三步：绝对连续 $\Rightarrow$ 几乎处处链式法则 $\Rightarrow$ 引理~\ref{lem:case-finite-time}。它不需要 Clarke 集值理论，也不需要预先验证原点不变性（引理已经包含），因此是最简洁、最值得推荐的路线。

\subsection{方法二：Dini 上导数}

Dini 上导数定义为（与比较引理~\ref{lem:comparison} 中用的导数一致）
\[
D^+V(t,x;f)=\limsup_{h\to0^+}\frac{V\big(x+hf(t,x)\big)-V(x)}{h}.
\]
它的优点是只依赖函数本身与\textbf{单点}向量场 $f$，不需要广义梯度。

\textbf{远离原点：} 此时 $f$ 单点，Dini 导数与普通导数一致，于是
\[
\text{(S1): }\ D^+V_2=-1,
\qquad
\text{(S2): }\ D^+V_2=-1+0.9\,\sgn{x}\sin t\le-0.1 .
\]

\textbf{原点处：} 若把 $f$ 换成整个 Filippov 集并取上确界，情况立刻变坏。对 $V_2=|x|$，$D^+V_2(0;v)=|v|$，故
\[
\sup_{v\in\mathbb{K}_2(t,0)}D^+V_2(0;v)
=\max\big\{|1+0.9\sin t|,\;|1-0.9\sin t|\big\}
=1+0.9|\sin t|>0 .
\]
也就是说，“对所有 Filippov 速度取最大”的 Dini 条件在原点处\textbf{必然失效}（$(S1)$ 给出 $1>0$，$(S2)$ 给出 $1+0.9|\sin t|>0$）。这不是例子不巧，而是方法本身的性质：原点处的允许速度包含“向外”的分量，而逐点取最大恰好挑中了它。

因此正确用法是：\textbf{远离原点}用 Dini 不等式（与普通导数一致），\textbf{原点}用引理~\ref{lem:case-finite-time} 的“到达后不离开”或 Dini 版 Lyapunov 定理（定理~\ref{thm:dini-lyapunov}，要求对每一个解分别验证）。不要指望一个全局的“取最大”条件。

\subsection{方法三：Clarke 广义梯度与集值 Lie 导数}

$V_2=|x|$ 的 Clarke 广义梯度（定理~\ref{thm:general-gradient}）为
\[
\partial V_2(x)=
\begin{cases}
  \ \{1\}, & x>0,\\
  \,[-1,1], & x=0,\\
  \ \{-1\}, & x<0 .
\end{cases}
\]

\textbf{远离原点：} 梯度是单点，$\zeta=\sgn{x}$，于是与前面一致：
\[
\zeta v=-1\ \ \text{((S1))},
\qquad
\zeta v\le-0.1\ \ \text{((S2))},
\qquad
\zeta\in\partial V_2(x),\ v\in\mathbb{K}_i(t,x).
\]

\textbf{原点处：} 令 $\xi\in\partial V_2(0)=[-1,1]$，$v\in\mathbb{K}_2(t,0)=0.9\sin t+[-1,1]$，按 Shevitz--Paden 的交集型集值 Lie 导数（定理~\ref{thm:set-valued-lie}）计算：
\[
\dot{\tilde V}_2(t,0)=\bigcap_{\xi\in[-1,1]}\xi^\top\mathbb{K}_2(t,0)=\{0\}.
\]
这不是偶然。一般地，只要 $0\in\mathbb{K}[f](x)$ 且 $\mathbb{K}[f](x)$ 有界，就有
\[
\bigcap_{\xi\in\partial V(x)}\xi^\top\mathbb{K}[f](x)=\{0\},
\]
因为 $0$ 属于每个集合 $\xi^\top\mathbb{K}[f](x)$，而当 $|\xi|\to0$ 时该集合收缩为 $\{0\}$。这个事实有两层含义：

\begin{itemize}
  \item 好的一面：$\dot{\tilde V}_2=\{0\}\subseteq(-\infty,0]$ 在原点自动成立，于是全局 $\dot{\tilde V}_2\subseteq(-\infty,0]$，$S=\{x:0\in\dot{\tilde V}_2(x)\}=\{0\}$，配合定理~\ref{thm:LaSalle}（需 Filippov 解唯一与紧不变集，本例都满足）可断言每条轨迹收敛到 $0$。
  \item 局限的一面：交集定义只保证“存在匹配的速度使内积为零”，并不保证“每个速度都使 $V$ 下降”，因此它在原点处给不出任何下降\textbf{速率}。有限时间结论仍然要靠第\ref{sec:case-methods}节方法一的 $\dot V_2\le-0.1$（$x\neq0$）与引理~\ref{lem:case-finite-time}。
\end{itemize}

作为对比，若采用“对所有 $(\zeta,v)$ 取最大”的 Clarke 条件，原点处会得到
\[
\max_{\zeta\in\partial V_2(0),\,v\in\mathbb{K}_2(t,0)}\zeta v=1+0.9|\sin t|>0,
\]
条件失败。所以三种方法都必须把原点单独处理——区别只在于：方法一用“到达后不离开”，方法二、三用“远离原点的不等式 + 不变性”。

\subsection{三种方法对照}

\begin{table}[!htb]
  \centering
  \small
  \begin{tabular}{@{}p{2.6cm}|p{3.6cm}|p{3.6cm}|p{3.4cm}@{}}
    \hline
    方法 & 需要的假设 & 原点处给出什么 & 结论强度 \\
    \hline
    几乎处处导数 & 解绝对连续；$V$ 沿解 a.e. 可导 & 不需要在原点求导：引理~\ref{lem:case-finite-time} 直接给出归零且不离开 & 有限时间 + 到达时间界（最强） \\
    \hline
    Dini 上导数 & $V$ 连续、局部 Lipschitz；单点 $f$ & “取最大”形式失效（$=1+0.9|\sin t|>0$），须逐解验证或分离处理 & 可达性、有限时间（需额外论证） \\
    \hline
    Clarke + 集值 Lie 导数 & $V$ 局部 Lipschitz、正则 & 交集自动为 $\{0\}$：弱不增，无速率 & 不变集原理（LaSalle）、渐近稳定 \\
    \hline
  \end{tabular}
  \caption{三种处理非光滑性的方法}
  \label{tab:case-methods}
\end{table}

三者的关系可以归纳为一句话：\textbf{都归结为“远离原点的严格下降 + 原点的强不变性”，差别只在“在原点附近如何合法地谈论导数”。}对 $(S1)$、$(S2)$ 这种一维系统，方法一最省事；方法二、三的价值在于能推广到多变量、非一致收敛、解不唯一的场合。

\section{Utkin 失效的二维反例}\label{sec:utkin-counterexample}

等效控制的四条条件中，任何一条被破坏都可能失效。前面已看到条件 3（可容许性）的破坏会导致间歇滑模；本节给出条件 2（横截条件）被破坏的反例，此时等效控制连定义都成问题。

\subsection{反例构造}

考虑如下二维系统：
\[
\dot{\boldsymbol x} = \begin{bmatrix} -x_2 \\ x_1 \end{bmatrix} + \begin{bmatrix} 1 \\ 1 \end{bmatrix} u,
\qquad
s(x) = x_1 - x_2,
\qquad
u = -k\,\sgn{s},\quad k>0.
\]
这里 $f(x) = \begin{bmatrix} -x_2 \\ x_1 \end{bmatrix}$，$g(x) = \begin{bmatrix} 1 \\ 1 \end{bmatrix}$。

\subsection{横截条件的破坏}

计算 $\nabla s = \begin{bmatrix} 1 \\ -1 \end{bmatrix}$，于是
\[
\nabla s^\top g = \begin{bmatrix} 1 & -1 \end{bmatrix} \begin{bmatrix} 1 \\ 1 \end{bmatrix} = 1 - 1 = 0.
\]
横截条件 $\nabla s^\top g\neq0$ 被破坏：控制 $u$ 的变化对 $\dot s$ 没有一阶影响。

\subsection{等效控制的失效}

按 Utkin 方法，令 $\dot s=0$：
\[
\dot s = \nabla s^\top (f + g u) = \begin{bmatrix} 1 & -1 \end{bmatrix} \left( \begin{bmatrix} -x_2 \\ x_1 \end{bmatrix} + \begin{bmatrix} 1 \\ 1 \end{bmatrix} u \right)
       = (-x_2 - x_1) + (1 - 1) u
       = -(x_1 + x_2).
\]
注意 $u$ 的系数为零，方程 $\dot s=0$ 简化为 $x_1+x_2=0$，与 $u$ 无关。这意味着：
\begin{itemize}
  \item 在 $s=0$ 与 $x_1+x_2=0$ 的交点（即原点 $x_1=x_2=0$）上，任何 $u$ 都能使 $\dot s=0$——\textbf{等效控制有无穷多个}；
  \item 在 $s=0$ 但 $x_1+x_2\neq0$ 的点上，$\dot s\neq0$ 且无法通过调节 $u$ 改变——\textbf{不存在任何等效控制}能维持滑模。
\end{itemize}
Utkin 的单点等效动力学 $u_\mathrm{eq} = -(\nabla s^\top g)^{-1}\nabla s^\top f$ 因分母为零而直接无定义。

\subsection{Filippov 框架下的良定描述}

Filippov 微分包含不受此影响：在 $s=0$ 上，$u=-k\,\sgn{s}$ 集值化为 $[-k,k]$，于是
\[
\mathbb{K}[f+gu](x) = f(x) + g(x)\,[-k,k]
                    = \begin{bmatrix} -x_2 \\ x_1 \end{bmatrix} + \begin{bmatrix} 1 \\ 1 \end{bmatrix} [-k,k]
                    = \begin{bmatrix} -x_2 + [-k,k] \\ x_1 + [-k,k] \end{bmatrix}.
\]
这是一个良定的集值映射，完整描述了滑模面上所有允许的速度；可以在此基础上讨论解的存在性与唯一性。

\begin{remark}[这个反例真正说明了什么]
  集合 $\mathbb{K}[f+gu]$ 良定，并不意味着“滑模仍然存在”。在 $s=0$ 上，任何允许速度 $v$ 的法向分量都是 $\nabla s^\top v=-(x_1+x_2)$，\textbf{与 $u$ 无关}。因此除原点外它恒不为零，轨迹总是横穿 $s=0$，\textbf{不存在非平凡的滑动模态}；在原点处法向分量恒为 $0$，$\mathbb{K}[f+gu](0)=\{(v,v):v\in[-k,k]\}$ 全部落在切向，滑模退化成一个平衡点。
  换句话说，Filippov 框架在这里的价值不是“救回”滑模，而是把“到底有没有滑动解、滑动速度是什么”变成\textbf{可以判定}的问题：Utkin 方法只能给出“$u_\mathrm{eq}$ 无定义”，Filippov 框架能给出“不存在非平凡滑动模态”这一完整结论，并且说明了原因是滑模面的法向方向不受控制影响，而不是控制增益不够大。
\end{remark}

\subsection{另一类失效场景：角点滑动}

除横截条件破坏外，还有一类 Utkin 方法失效的典型场景：\textbf{多张滑模面相交的角点}。当系统有多个切换面 $\{s_1=0\},\{s_2=0\},\dots$ 且在某个角点（如 $s_1=s_2=0$）相交时，Utkin 的单滑模面等效控制方法无法直接推广；Filippov 微分包含可以自然地处理这种多个集值集合交叠的情形，给出角点处的滑动解定义，而 Utkin 范式需要额外的几何条件与分情况讨论。

\section{核心对比：Utkin vs.\ Filippov}\label{sec:comparison}

\begin{table}[!htp]
  \centering
  \small
  \begin{tabular}{|p{6.5em}|p{16.8em}|p{16.8em}|}
    \hline
    \textbf{对比项} & \textbf{Utkin 等效控制} & \textbf{Filippov 微分包含 + Clarke 次梯度 / 集值李导数} \\
    \hline
    定位 & 工程设计工具 & 底层数学框架 \\
    \hline
    核心操作 & 拆分趋近 / 滑动阶段，反解等效控制 $u_\mathrm{eq}$ & 将不连续右端集值化，用微分包含描述所有允许解 \\
    \hline
    \makecell[l]{Lyapunov\\工具} & 光滑 $V=\tfrac12 s^2$，普通李导数 & 非光滑 $V$，Clarke 广义梯度 $\partial V$，集值李导数 $\dot{\tilde V}$ 或 Dini 上导数 \\
    \hline
    适用条件 & 仿射系统、单张光滑滑模面、$\nabla s^\top g\neq0$、$|u_\mathrm{eq}|\le k$、可达性 & 只需右端可测、局部本质有界 \\
    \hline
    滑模面上 & 单点等效动力学 $\dot x = f + g u_\mathrm{eq}$（降维连续系统） & 集值速度集合 $\mathbb{K}[f+gu](x)$，可能包含多个滑动方向 \\
    \hline
    适用场景 & 常规单滑模面仿射系统 & 单滑模面、多滑模面相交、横截条件破坏、奇异控制等复杂情形 \\
    \hline
    主要用途 & 控制器设计、仿真、物理解释 & 解的存在唯一性严格证明、不变集原理、非光滑 Lyapunov 稳定性分析 \\
    \hline
  \end{tabular}
  \caption{Utkin 等效控制与 Filippov 微分包含的对比}
  \label{tab:utkin-vs-filippov}
\end{table}

\section{本章小结}\label{sec:case-summary}

把两条系统的结论汇总如下（均指 Filippov 意义）：

\begin{table}[!htb]
  \centering
  \small
  \begin{tabular}{@{}p{5cm}|c|c@{}}
    \hline
    性质 & (S1) $\dot x=-\sgn{x}$ & (S2) $\dot x=-\sgn{x}+0.9\sin t$ \\
    \hline
    原点是否为平衡点 / Filippov 解 & 是 & 是 \\
    \hline
    原点的强不变性 & 是 & 是 \\
    \hline
    全局 Lyapunov 稳定 & 是 & 是 \\
    \hline
    全局有限时间稳定 & 是 & 是 \\
    \hline
    到达时间界 & $|x_0|$（精确） & $10\,|x_0|$ \\
    \hline
    等效控制 $u_\mathrm{eq}$ & $0$ & $-0.9\sin t$（可容许） \\
    \hline
    Carathéodory 全局解 & 有（靠 $\sgn{0}=0$） & 无 \\
    \hline
  \end{tabular}
  \caption{两个系统的结论对照}
  \label{tab:case-conclusions}
\end{table}

\textbf{推荐做法：}取 $V=|x|$，沿 Filippov 解使用几乎处处导数得到 $\dot V\le-1$ 或 $\dot V\le-0.1$（$x\neq0$），再用引理~\ref{lem:case-finite-time} 一次性得到“有限时间到达 + 到达后不离开”。这条路线只需三步，且对两条系统统一适用。

最后回到全书的定位。两套范式并非对立，而是\textbf{互补}关系：

\begin{itemize}
  \item \textbf{Utkin 等效控制}偏向工程设计，简单好用但假设强。在仿射结构、单张光滑滑模面、横截条件与可容许条件同时成立时，它给出的等效动力学与 Filippov 框架严格等价（命题~\ref{prop:case-ueq}），是滑模控制教材与工程实践的主流工具。
  \item \textbf{Filippov 微分包含 + 非光滑 Lyapunov 工具}是底层数学框架，适用范围更广。它不依赖仿射结构与横截条件，能统一处理滑模面相交、横截条件破坏、间歇滑模等奇异情形，并为解的存在性、不变集原理与非光滑稳定性提供严格证明。
  \item 二者的分工可以用本章的一句话概括：\textbf{等效控制负责滑模面，非光滑工具负责滑模面之外}。教材教你“怎么用”，论文回答“为什么能用”。
\end{itemize}

第\ref{ch:lipschitz}章回答“经典解什么时候存在且唯一”，第\ref{ch:bounded}章回答“右端不连续时解与稳定性如何定义”，本章回答“在具体的滑模控制问题里应该选哪一套、各自的前提是什么”。更接近算法层面的内容（超螺旋算法的增益条件与严格 Lyapunov 证明）见第\ref{ch:sta}章。
